% Figure: drying rate curve, constant-rate then falling-rate period vs moisture content.
\begin{figure}[htbp]
\centering
\begin{tikzpicture}
\begin{axis}[
  width=11cm, height=6.5cm,
  xlabel={free moisture content $X$ (kg water / kg dry solid)},
  ylabel={drying rate $R$},
  xmin=0, xmax=1.0, ymin=0, ymax=1.1,
  ytick=\empty,
  legend pos=north west,
  legend cell align=left,
  grid=both, grid style={enggray!18},
  every axis plot/.append style={very thick},
]
% constant-rate period from Xc=0.4 up to X=0.9 (rate = 1.0)
\addplot[engblue] coordinates {(0.4,1.0) (0.9,1.0)};
\addlegendentry{constant-rate period}
% falling-rate period: linear from Xc down to X=0 (rate from 1.0 to 0)
\addplot[engred] coordinates {(0.0,0.0) (0.4,1.0)};
\addlegendentry{falling-rate period}
% critical moisture marker
\draw[enggray, dashed] (axis cs:0.4,0) -- (axis cs:0.4,1.0);
\node[font=\scriptsize, anchor=south] at (axis cs:0.4,1.02) {$X_c$};
\node[font=\scriptsize, anchor=west] at (axis cs:0.42,0.02) {critical moisture};
\end{axis}
\end{tikzpicture}
\caption[Drying-rate curve]{The drying-rate curve. Above the critical
moisture content $X_c$ the solid surface stays wetted and the rate is set
by air-side mass transfer at the saturation deficit, holding nearly
constant. Below $X_c$ liquid can no longer reach the surface as fast as
it would evaporate, so the rate falls toward zero as the solid dries out.
The constant-rate period clears the bulk free water; the long falling-rate
tail removes the last bound moisture and usually dominates the total
drying time.}
\label{fig:vol04ch07:drying-curve}
\end{figure}
